% 復習4　文献を批判的に読む
\session{4}{1}{90分}{文献を批判的に読む}{AI要約と対話}{}

\goals{
  \item 論文の構造（問い、方法、結果、考察、限界）を押さえて読める
  \item AI要約の使いどころと限界を説明できる。要約は原文の代わりにならない
  \item 主張と根拠の対応、データの範囲、別の解釈の可能性を確かめる批判的な読み方ができる
}

\fillin{読む文献（著者・年・題名）}

\begin{caution}{入力するときの注意}
論文の全文を入力する場合は、大学のAI利用規程と、その資料の利用条件を確認する。要旨や必要な部分だけを使えば十分なことが多い。
\end{caution}

\work[80mm]{ワーク1}{AI要約と原文を照らし合わせる}
\begin{promptbox}[構造をつかむ]
次の論文について、「問い」「方法」「主な結果」「著者が認めている限界」をそれぞれ2文以内でまとめてください。書かれていない項目は「記載なし」としてください。（論文の要旨や該当部分を貼る）
\end{promptbox}
\begin{wstabh}{colspec={Q[l,wd=17mm,m] X[l,m] X[l,m] X[l,m]},row{2-Z}={ht=17mm}}
  & AIの要約（要点） & 原文で確認した内容（頁） & 要約の誤り・抜け \\
  問い & & & \\ 方法 & & & \\ 主な結果 & & & \\ 限界 & & & \\
\end{wstabh}

\work[70mm]{ワーク2}{最も強い反論に、原文の根拠で答える}
\begin{promptbox}[反論を引き出す]
この論文の主張「（主張）」に対して、最も強い反論を3つ挙げてください。それぞれについて、どのようなデータがあれば検証できるかも示してください。
\end{promptbox}
\begin{wstabh}{colspec={X[3,l,m] Q[c,wd=15mm,m] X[3,l,m]},row{2-Z}={ht=16mm}}
  AIが挙げた反論 & 答えられる\newline ○△× & 原文の根拠（頁）・答えられない理由 \\
  & & \\ & & \\ & & \\
\end{wstabh}

\work{ワーク3}{未解決の要素}
\twoboxes{論文が答えていないこと}{自分の問いとの隙間}{32mm}

\begin{nexttime}
  \item 読んだ文献について「主張」「根拠」「限界」を各100字程度でまとめる
  \item AIとの議論で見えた未解決の要素を1つ書き、自分の問いとどうつながるかを考える
\end{nexttime}
\answerbox[主張（100字程度）]{24mm}
\answerbox[根拠（100字程度）]{24mm}
\answerbox[限界（100字程度）]{24mm}
\answerbox[未解決の要素と、自分の問いとのつながり]{24mm}
\aimemo
